The first complete thermal runaway reaction network implemented in
LiionTR is based on the reduced-order model reported by
\textcite{hu2020inhibition}.

The model represents thermal runaway using four principal reaction
states:

\begin{enumerate}
    \item solid-electrolyte interphase (SEI) decomposition;
    \item anode-electrolyte reaction;
    \item cathode decomposition;
    \item electrolyte decomposition.
\end{enumerate}

Rather than implementing the model as a single monolithic component,
LiionTR represents each physical contribution using the generic kinetic,
progress, reaction, and reaction-network abstractions described in
\cref{sec:kinetics,sec:reactions}.


\subsection{Reaction-state vector}

The reaction-state vector is

\begin{equation}
    \boldsymbol{\alpha}
    =
    \begin{bmatrix}
        \alpha_{\mathrm{SEI}} \\
        \alpha_{\mathrm{an}} \\
        \alpha_{\mathrm{cat}} \\
        \alpha_{\mathrm{el}}
    \end{bmatrix},
    \label{eq:hu_state_vector}
\end{equation}

where the subscripts denote SEI decomposition, anode-electrolyte
reaction, cathode decomposition, and electrolyte decomposition,
respectively.

The initial state used by the current LiionTR implementation is

\begin{equation}
    \boldsymbol{\alpha}_0
    =
    \begin{bmatrix}
        0.85 \\
        0.25 \\
        0.04 \\
        0.00
    \end{bmatrix}.
    \label{eq:hu_initial_state}
\end{equation}


\subsection{Arrhenius formulation}

The thermally activated reactions use the Arrhenius expression

\begin{equation}
    k_i(T)
    =
    A_i
    \exp
    \left(
        -\frac{E_{a,i}}{RT}
    \right),
    \label{eq:hu_arrhenius}
\end{equation}

where $A_i$ is the pre-exponential factor and $E_{a,i}$ is the
activation energy of reaction $i$.

The resulting conversion rate is obtained by multiplying the kinetic
term by the appropriate reaction-progress model,

\begin{equation}
    \dot{\alpha}_i
    =
    k_i(T)
    f_i
    \left(
        \alpha_i,
        \boldsymbol{\alpha}
    \right).
    \label{eq:hu_generic_progress}
\end{equation}


\subsection{SEI decomposition}

The SEI decomposition reaction is represented using first-order
remaining-fraction kinetics,

\begin{equation}
    \frac{d\alpha_{\mathrm{SEI}}}{dt}
    =
    k_{\mathrm{SEI}}(T)
    \left(
        1-\alpha_{\mathrm{SEI}}
    \right).
    \label{eq:hu_sei_progress}
\end{equation}

The implemented kinetic parameters are

\begin{align}
    A_{\mathrm{SEI}}
    &=
    \num{1.667e15}
    \ \si{\per\second},
    \\
    E_{a,\mathrm{SEI}}
    &=
    \num{1.3508e5}
    \ \si{\joule\per\mole},
    \\
    \Delta H_{\mathrm{SEI}}
    &=
    \num{2.57e5}
    \ \si{\joule\per\kilogram}.
\end{align}

The reported volumetric reactant content is converted to SI units as

\begin{equation}
    \num{6.104e5}
    \ \si{\gram\per\meter\cubed}
    =
    \num{610.4}
    \ \si{\kilogram\per\meter\cubed}.
\end{equation}

The literature model specifies an initial remaining fraction of 0.15.
Because LiionTR uses conversion,

\begin{equation}
    \alpha_{\mathrm{SEI},0}
    =
    1-0.15
    =
    0.85.
\end{equation}


\subsection{Anode-electrolyte reaction}

The anode-electrolyte reaction uses the Arrhenius parameters

\begin{align}
    A_{\mathrm{an}}
    &=
    \num{2.5e13}
    \ \si{\per\second},
    \\
    E_{a,\mathrm{an}}
    &=
    \num{1.3508e5}
    \ \si{\joule\per\mole},
    \\
    \Delta H_{\mathrm{an}}
    &=
    \num{1.714e6}
    \ \si{\joule\per\kilogram}.
\end{align}

The reaction enthalpy follows from the unit conversion

\begin{equation}
    \SI{1714}{\joule\per\gram}
    =
    \num{1.714e6}
    \ \si{\joule\per\kilogram}.
\end{equation}

The same volumetric reactant content as the SEI reaction is used,

\begin{equation}
    W_{\mathrm{an}}
    =
    \SI{610.4}{\kilogram\per\meter\cubed}.
\end{equation}

The initial remaining fraction is 0.75, corresponding to

\begin{equation}
    \alpha_{\mathrm{an},0}
    =
    1-0.75
    =
    0.25.
\end{equation}


\subsubsection{Coupling to SEI decomposition}

The anode-electrolyte reaction is not represented as an independent
first-order process. Its progress depends on the state of the SEI
reaction.

The remaining SEI fraction is

\begin{equation}
    r_{\mathrm{SEI}}
    =
    1-\alpha_{\mathrm{SEI}}.
\end{equation}

In the current LiionTR implementation, the anode reaction becomes active
only when

\begin{equation}
    r_{\mathrm{SEI}}
    <
    0.10.
    \label{eq:hu_anode_sei_threshold}
\end{equation}

The reaction is additionally modified by an exponential inhibition term
derived from the normalized SEI state.

The resulting model may be written conceptually as

\begin{equation}
    \frac{d\alpha_{\mathrm{an}}}{dt}
    =
    k_{\mathrm{an}}(T)
    \left(
        1-\alpha_{\mathrm{an}}
    \right)
    I_{\mathrm{SEI}},
    \label{eq:hu_anode_progress}
\end{equation}

where

\begin{equation}
    I_{\mathrm{SEI}}
    =
    \exp
    \left(
        -x_{\mathrm{SEI}}
    \right)
\end{equation}

when the SEI threshold condition in
\cref{eq:hu_anode_sei_threshold} is satisfied; otherwise the reaction
progress rate is zero.

This formulation illustrates the context-dependent reaction architecture
implemented in LiionTR.


\subsection{Cathode decomposition}

Cathode decomposition differs from the preceding reactions because it is
represented using two parallel kinetic channels sharing a single
reaction conversion,

\begin{equation}
    \alpha_{\mathrm{cat}}.
\end{equation}

The cathode reaction is activated only above

\begin{equation}
    T_{\mathrm{cat}}
    =
    \SI{393.15}{\kelvin}.
    \label{eq:hu_cathode_threshold}
\end{equation}

The reported cathode content is

\begin{equation}
    \num{1.221e6}
    \ \si{\gram\per\meter\cubed},
\end{equation}

which LiionTR converts to

\begin{equation}
    W_{\mathrm{cat}}
    =
    \SI{1221}{\kilogram\per\meter\cubed}.
\end{equation}


\subsubsection{Cathode channel 1}

The first kinetic channel uses

\begin{align}
    A_{\mathrm{cat},1}
    &=
    \num{1.75e9}
    \ \si{\per\second},
    \\
    E_{a,\mathrm{cat},1}
    &=
    \num{1.1495e5}
    \ \si{\joule\per\mole},
    \\
    \Delta H_{\mathrm{cat},1}
    &=
    \num{7.7e4}
    \ \si{\joule\per\kilogram}.
\end{align}


\subsubsection{Cathode channel 2}

The second kinetic channel uses

\begin{align}
    A_{\mathrm{cat},2}
    &=
    \num{1.077e12}
    \ \si{\per\second},
    \\
    E_{a,\mathrm{cat},2}
    &=
    \num{1.5888e5}
    \ \si{\joule\per\mole},
    \\
    \Delta H_{\mathrm{cat},2}
    &=
    \num{8.4e4}
    \ \si{\joule\per\kilogram}.
\end{align}


\subsubsection{Shared cathode progress}

The cathode reaction uses the autocatalytic progress function

\begin{equation}
    f_{\mathrm{cat}}
    =
    \alpha_{\mathrm{cat}}
    \left(
        1-\alpha_{\mathrm{cat}}
    \right).
    \label{eq:hu_cathode_progress}
\end{equation}

The shared cathode conversion therefore evolves according to

\begin{equation}
    \frac{d\alpha_{\mathrm{cat}}}{dt}
    =
    f_{\mathrm{cat}}
    \left[
        k_{\mathrm{cat},1}(T)
        +
        k_{\mathrm{cat},2}(T)
    \right].
    \label{eq:hu_cathode_conversion}
\end{equation}

The initial conversion is

\begin{equation}
    \alpha_{\mathrm{cat},0}
    =
    0.04.
\end{equation}

Both kinetic channels use the same progress variable, and their heat
contributions are combined.

The resulting mass-specific heat-generation rate is

\begin{equation}
    \dot{q}_{\mathrm{cat}}
    =
    w_{\mathrm{cat}}
    f_{\mathrm{cat}}
    \left[
        \Delta H_{\mathrm{cat},1}
        k_{\mathrm{cat},1}(T)
        +
        \Delta H_{\mathrm{cat},2}
        k_{\mathrm{cat},2}(T)
    \right],
    \label{eq:hu_cathode_heat}
\end{equation}

for temperatures satisfying the threshold condition in
\cref{eq:hu_cathode_threshold}.


\subsection{Electrolyte decomposition}

Electrolyte decomposition is represented by an Arrhenius reaction
activated above

\begin{equation}
    T_{\mathrm{el}}
    =
    \SI{473.15}{\kelvin}.
    \label{eq:hu_electrolyte_threshold}
\end{equation}

The kinetic and energetic parameters are

\begin{align}
    A_{\mathrm{el}}
    &=
    \num{5.14e25}
    \ \si{\per\second},
    \\
    E_{a,\mathrm{el}}
    &=
    \num{2.74e5}
    \ \si{\joule\per\mole},
    \\
    \Delta H_{\mathrm{el}}
    &=
    \num{1.55e5}
    \ \si{\joule\per\kilogram}.
\end{align}

The reported reaction enthalpy is converted according to

\begin{equation}
    \SI{155}{\joule\per\gram}
    =
    \num{1.55e5}
    \ \si{\joule\per\kilogram}.
\end{equation}

The reported volumetric electrolyte content is

\begin{equation}
    \num{4.069e5}
    \ \si{\gram\per\meter\cubed},
\end{equation}

which becomes

\begin{equation}
    W_{\mathrm{el}}
    =
    \SI{406.9}{\kilogram\per\meter\cubed}.
\end{equation}

The initial remaining fraction is unity, giving

\begin{equation}
    \alpha_{\mathrm{el},0}
    =
    0.
\end{equation}


\subsection{Summary of model parameters}

The principal parameters used by the current LiionTR implementation are
summarized in \cref{tab:hu_parameters}.

\begin{table}[htbp]
    \centering
    \small

    \caption{
        Principal kinetic and energetic parameters used in the LiionTR
        implementation of the thermal runaway model of
        \textcite{hu2020inhibition}.
    }
    \label{tab:hu_parameters}

    \begin{tabular}{
        l
        S[scientific-notation=true]
        S[scientific-notation=true]
        S[scientific-notation=true]
        S
    }
        \toprule
        {Reaction}
        & {$A$ (\si{\per\second})}
        & {$E_a$ (\si{\joule\per\mole})}
        & {$\Delta H$ (\si{\joule\per\kilogram})}
        & {$\alpha_0$}
        \\
        \midrule

        SEI
        & 1.667e15
        & 1.3508e5
        & 2.57e5
        & 0.85
        \\

        Anode-electrolyte
        & 2.5e13
        & 1.3508e5
        & 1.714e6
        & 0.25
        \\

        Cathode channel 1
        & 1.75e9
        & 1.1495e5
        & 7.7e4
        & 0.04
        \\

        Cathode channel 2
        & 1.077e12
        & 1.5888e5
        & 8.4e4
        & 0.04
        \\

        Electrolyte
        & 5.14e25
        & 2.74e5
        & 1.55e5
        & 0.00
        \\

        \bottomrule
    \end{tabular}
\end{table}


\subsection{Volumetric content and effective mass fractions}

Hu et al.\ report several reacting-material inventories on a volumetric
basis \parencite{hu2020inhibition}.

LiionTR converts these quantities to effective mass fractions using

\begin{equation}
    w_i
    =
    \frac{
        W_i
    }{
        \rho_{\mathrm{cell}}
    },
    \label{eq:hu_mass_fraction}
\end{equation}

where $W_i$ is the volumetric reactant content and
$\rho_{\mathrm{cell}}$ is the effective cell density.

The total heat generation from reaction $i$ is consequently

\begin{equation}
    \dot{Q}_i
    =
    m_{\mathrm{cell}}
    w_i
    \Delta H_i
    \dot{\alpha}_i.
    \label{eq:hu_reaction_heat}
\end{equation}


\subsection{Complete thermal coupling}

The total heat-generation rate is

\begin{equation}
    \dot{Q}_{\mathrm{gen}}
    =
    \dot{Q}_{\mathrm{SEI}}
    +
    \dot{Q}_{\mathrm{an}}
    +
    \dot{Q}_{\mathrm{cat}}
    +
    \dot{Q}_{\mathrm{el}}.
    \label{eq:hu_total_heat}
\end{equation}

This source term is coupled directly to the lumped thermal balance,

\begin{equation}
    C_{\mathrm{th}}
    \frac{dT}{dt}
    =
    \dot{Q}_{\mathrm{gen}}
    -
    hA
    \left(
        T-T_{\infty}
    \right).
    \label{eq:hu_thermal_balance}
\end{equation}

The Arrhenius temperature dependence then creates the thermal feedback
responsible for runaway behavior.


\subsection{LiionTR representation}

The mapping of the Hu et al.\ model onto the LiionTR architecture is
illustrated in \cref{fig:hu_architecture}.

\begin{figure}[htbp]
    \centering

    \begin{tikzpicture}[
        node distance=1.35cm and 2.0cm,
        block/.style={
            rectangle,
            rounded corners,
            draw,
            align=center,
            minimum width=3.4cm,
            minimum height=0.9cm
        },
        line/.style={
            -{Latex[length=3mm]},
            thick
        }
    ]

    \node[block] (sei) {SEI decomposition};

    \node[block, below left=of sei] (anode)
        {Anode-electrolyte};

    \node[block, below right=of sei] (cathode)
        {Cathode decomposition\\two kinetic channels};

    \node[block, below=of $(anode)!0.5!(cathode)$] (electrolyte)
        {Electrolyte decomposition};

    \node[block, below=of electrolyte] (network)
        {Reaction network};

    \node[block, below left=of network] (rates)
        {Reaction progress rates};

    \node[block, below right=of network] (heat)
        {Total heat generation};

    \draw[line] (sei) -- (anode);
    \draw[line] (sei) -- (network);
    \draw[line] (anode) -- (network);
    \draw[line] (cathode) -- (network);
    \draw[line] (electrolyte) -- (network);

    \draw[line] (network) -- (rates);
    \draw[line] (network) -- (heat);

    \end{tikzpicture}

    \caption{
        Representation of the Hu et al.\ thermal runaway model within
        the generic LiionTR reaction-network architecture.
    }
    \label{fig:hu_architecture}
\end{figure}

The arrow between the SEI and anode blocks indicates that the
anode-electrolyte reaction depends on the state of the SEI reaction.
The remaining reactions are assembled into the same network and coupled
through the common thermal state.


\subsection{Scope of the implementation}

The present implementation focuses on the reduced-order reaction and
heat-generation structure described by \textcite{hu2020inhibition}.

The LiionTR representation should therefore be distinguished from a
complete mechanistic electrochemical or thermochemical model.

In particular, the current Hu reaction network does not itself determine

\begin{itemize}
    \item equilibrium gas composition;
    \item gas-phase secondary reactions;
    \item electrochemical electrode potentials;
    \item spatial temperature distributions;
    \item mechanical deformation or rupture;
    \item multiphase vent products.
\end{itemize}

These phenomena are intended to be represented by separate LiionTR
subsystems or external scientific backends.